\section{Discussion}
\label{sec:discussion}

The central finding of this paper is that the Tiebout sorting mechanism is remarkably resilient to parasitism. When extremist communities capable of extracting utility from mainstream neighbors enter a diverse platform ecosystem, the sorting process absorbs the threat without catastrophic welfare loss: the full range of normalized mainstream utility across the entire factorial spans only 0.135, a margin that would be difficult to distinguish from noise in most empirical settings. Communities displaced by extremist raids typically land on better-matched platforms, and the system self-corrects after each disruption. This resilience, however, is not unconditional. The model identifies a specific structural failure mode---the intersection of high parasitism intensity, direct-voting governance, and sufficient platform diversity for extremists to form stable exploitation equilibria---under which a subset of platforms absorbs disproportionate harm while the rest of the system appears healthy. The policy-relevant story is therefore not about average damage but about the conditions under which structural protection breaks down and the communities least equipped to defend themselves bear concentrated costs.

\subsection{Platform Design and the Consolidation Problem}

The diversification premium---the welfare gain from expanding the number of available platforms---is one of the most robust findings in the simulation. More platforms produce better outcomes, lower variance, and finer preference matching, consistent with the core Tiebout prediction. What the extremism experiments add is that this premium is not merely additive: it grows with the severity of the threat. The premium nearly doubles from low to high parasitism intensity, and the mechanism shifts from preference matching (which dominates in benign environments) to target dilution (which dominates under threat). More platforms spread extremist communities across more jurisdictions, reducing parasitism density on any given platform and preserving the sorting mechanism's ability to match mainstream communities with compatible governance.

This finding speaks directly to the contemporary platform landscape. The social media ecosystem is not characterized by the jurisdictional diversity that the Tiebout mechanism requires. A small number of firms---Meta, TikTok, Google---command the vast majority of user attention, and market dynamics continue to push toward consolidation rather than diversification. Scholars of media systems have long argued that viewpoint diversity depends on structural diversity in the media landscape \citep{tufekci2017twitter}; the model adds a sharper claim. Platform consolidation does not merely reduce preference matching---it systematically worsens outcomes for mainstream communities, and the magnitude of the harm scales with the extremist threat present. In a benign environment, consolidation leaves modest utility on the table. In a hostile one, it erodes the structural buffer that protects mainstream communities from parasitic exploitation. The tension is structural: the Tiebout mechanism requires a minimum of jurisdictional diversity to function, but the network effects and economies of scale that characterize platform markets push relentlessly toward consolidation.

Two qualifications temper this implication. First, the model's effects are modest on average. Consolidation alone does not produce catastrophe; it removes the system's capacity to absorb shocks when they arrive. The case for platform diversity is therefore less about preventing everyday harm and more about preserving structural resilience against tail-risk events---the platform equivalent of maintaining biodiversity not for its average-day benefits but for its insurance value. Second, the model treats platforms as passive and homogeneous within governance type. Real platforms differ in scale, network effects, content niche, and strategic behavior in ways that complicate a simple ``more is better'' prescription. The theoretical prediction is that jurisdictional diversity is protective; the policy translation requires attention to what kinds of new platforms would actually absorb communities and dilute extremist density, rather than simply fragmenting audiences across functionally identical services.

\subsection{Governance Design and the Coalition Paradox}

Coalition-based governance emerges as the strongest structural firewall against parasitic exploitation, and it does so through a mechanism that is genuinely counterintuitive. In the baseline experiments without extremists, coalition platforms appeared mediocre---slower to converge, lower in average utility than algorithmic platforms, and attracting a stable but unremarkable share of communities. Only when extremists entered the system did the coalition mechanism reveal its protective value. Mainstream utility on coalition platforms declined by just 0.22 units across the full range of parasitism intensity, compared to a decline of 3.60 on direct-voting platforms. The mechanism is enclave formation: coalition governance induces communities to form ideological blocs that jointly control policy dimensions, and this internal structure naturally segregates extremist and mainstream communities into separate coalitions. Once segregated, the parasitism channel is structurally severed---extremists and mainstream communities are no longer effective neighbors even when they share the same platform.

This enclave mechanism provides a structural analog to what the moderation literature describes as community self-governance. Platforms that support movement-formation tools---coalition voting, co-moderation, community-created rules---enable communities to construct internal boundaries that insulate members from hostile neighbors. The finding resonates with empirical observations from platforms like Reddit and Discord, where subreddit and server structures allow communities to establish and enforce their own norms, and where the most stable communities are those with the strongest internal governance \citep{chandrasekharan2017you}. The model suggests that the protective value of these structures derives not from their moderation rules per se but from the coalition-formation process itself: the act of aggregating preferences within blocs produces ideological coherence that extremists cannot easily penetrate.

The practical implication is that governance mechanisms supporting coalition formation---community voting systems, co-moderation structures, movement-building tools---may be more protective than either algorithmic sorting or direct majority rule. Algorithmic platforms serve as effective refuges in mixed systems, absorbing the largest share of communities and delivering acceptable utility to heterogeneous preferences. But they do not actively protect against parasitism; they merely dilute it. Direct-voting platforms, by contrast, shift from boutique to vulnerable as parasitism intensifies: the same governance mechanism that rewarded preference-aligned communities at low threat levels exposes unprotected mainstream communities to exploitation at high ones. Coalition governance occupies a distinct structural position---it actively generates protection through endogenous segregation, and this protection strengthens rather than weakens as the threat grows.

A significant limitation applies here. The model assumes that community coalitions can form organically---that communities possess the information, motivation, and institutional support to identify compatible partners and negotiate joint control of policy dimensions. Real platforms vary enormously in how much they support this process. Many mainstream platforms offer no coalition-formation mechanisms at all, relying entirely on algorithmic personalization or platform-level policy. The finding therefore identifies a design principle---governance that supports community-level coalition formation is structurally protective---rather than a description of existing platform behavior.

\subsection{Extremist Strategy as a System Property}

The empirical puzzle that motivated this paper---why extremist movements cycle across platforms in the specific pattern of concentration, raiding, retreat, and re-entry---receives a structural explanation from the simulation results. The raiding pattern is an emergent property of the sorting process under mixed governance, not the product of strategic planning by any individual community. No extremist community in the model plans a raid cycle; the cycle emerges from individual utility-maximizing moves given the payoff structure of the platform ecosystem. Extremists accumulate on direct-voting platforms where unprotected mainstream neighbors offer parasitic payoffs, gradually shift platform policy toward their preferences, drive mainstream communities out, experience declining parasitic returns, and then disperse to platforms where new targets are available. The cycle repeats because the structural conditions that produced it persist.

The simulation results further reveal that parasitism intensity does not merely scale the magnitude of extremist harm---it produces qualitatively distinct behavioral regimes. At low intensity, extremists behave as ideologues, seeking platforms whose governance bundles align with their polarized preferences. At high intensity, the same communities become griefers, concentrating wherever mainstream targets are most accessible regardless of policy alignment. This regime shift means that the character of the extremist threat, not just its magnitude, changes with the structural conditions of the ecosystem.

This finding reframes how we interpret empirically observed cross-platform coordination. Episodes like the Kiwi Farms harassment campaigns, GamerGate-era cross-platform mobilization, and recurring patterns of extremist return to mainstream spaces after deplatforming look planned, and some coordination certainly occurs. But the model demonstrates that much of this cycling behavior can arise from purely structural incentives: given a platform ecosystem containing direct-voting platforms with unprotected mainstream neighbors, the raiding pattern is the equilibrium outcome of decentralized utility-maximizing behavior. The distinction matters for policy. If cycling is purely strategic, then disrupting coordination channels---encrypted messaging, off-platform planning spaces---should break the cycle. If cycling is substantially structural, then disrupting channels alone is insufficient; the cycle will re-emerge as long as the ecosystem contains the governance configurations that sustain it.

The displacement paradox reinforces this structural interpretation. When raids displace mainstream communities from direct platforms, those communities typically find better matches elsewhere---mean displacement utility is positive, and at high platform diversity no displacement events produce negative utility changes. The Tiebout mechanism self-corrects: raids are costly for the platforms that experience them but beneficial, in net, for the displaced communities and the system as a whole. This self-correction is itself a form of resilience, and it depends on the same jurisdictional diversity identified in the diversification results. At low platform counts, where sorting options are scarce, displacement occasionally produces negative outcomes. The system's capacity to absorb raiding shocks is a function of its structural diversity, closing the loop between the consolidation and extremist-strategy findings.

The policy implication is that deplatforming from a single platform does not break the extremist cycle if the structural conditions---direct-voting platforms with unprotected mainstream neighbors, sufficient platform diversity for re-entry---persist elsewhere in the ecosystem. Effective intervention requires attention to the architecture of the platform system, not only to the moderation decisions of individual platforms.

\subsection{Limitations}

Several modeling choices constrain the scope of these findings. First, community preferences are fixed throughout each simulation. Real communities update their preferences in response to platform experiences, including exposure to extremist content, and preference dynamics could amplify or dampen the sorting patterns observed here. Second, platforms are passive actors: governance type is assigned and does not change. Real platforms strategically adjust their moderation, algorithms, and features to attract or repel communities, and these adjustments would introduce feedback loops absent from the current model. Third, the three governance types---direct voting, coalition formation, and algorithmic recommendation---are stylizations of a much more complex reality. Most real platforms combine elements of all three, and the boundaries between types shift over time as platforms redesign their systems.

Fourth, while the experimental design specifies planned sensitivity analyses, those results are not yet reported. The robustness of the core findings to alternative parameterizations, utility functional forms, and movement rules remains to be confirmed. Fifth, the parameter ranges for platform count, extremist proportion, and parasitism intensity are theoretically motivated but not calibrated to specific real-world platforms or communities. The model is a theoretical device for identifying structural mechanisms, not a predictive tool for forecasting outcomes in particular platform ecosystems. These limitations are characteristic of the exploratory modeling tradition in which this work sits \citep{bankesExploratoryModelingPolicy1993}: the aim is to build theory and identify mechanisms that can guide empirical inquiry, not to generate point predictions.

\subsection{Conclusion}

This paper began with an empirical puzzle: why do extremist communities cycle across the platform ecosystem in the specific pattern of concentration, raiding, retreat, and re-entry? The deplatforming literature treats platform access as binary, and the moderation literature focuses on within-platform governance. Neither framework explains the system-level dynamics that make cycling both rational and emergent. The modified Tiebout model developed here provides that explanation. Community sorting across a system of heterogeneous platforms produces an equilibrium in which extremists endogenously concentrate on direct-voting platforms, exploit unprotected mainstream neighbors, and disperse when parasitic returns decline---generating the cycling pattern as a structural property of the ecosystem rather than a product of coordinated strategy.

The theoretical contribution is to extend the Tiebout framework to multi-governance platform ecosystems and to identify the structural conditions under which parasitic exploitation becomes self-sustaining. The policy contribution is to redirect attention from individual moderation decisions to the architecture of the platform system itself. It is the mix of governance types, the degree of jurisdictional diversity, and the availability of coalition-formation mechanisms that determine whether mainstream communities are structurally protected or structurally exposed. Future work should pursue empirical calibration against observed platform migration data, introduce dynamic governance that allows platforms to respond strategically to community composition, and model preference updating to capture the feedback between extremist exposure and community radicalization.
